\documentclass[12pt]{article}
\usepackage[utf8]{inputenc}
\usepackage[T1]{fontenc}
\usepackage{hyperref}
\title{Puquina Language Resources for K-12 Curriculum Development}
\author{Compiled by Claude Sonnet 5.5 (claude-sonnet-5-5) for Andr\'as Kornai}
\date{}
\begin{document}
\maketitle

\section{Introduction}
Puquina (ISO 639-3 \texttt{puq}, Glottocode puqu1242) is an extinct language once spoken around Lake Titicaca, between Arequipa and Tacna, and between Sucre and Potos\'i (Peru, Bolivia, Chile). It disappeared from the record in the early nineteenth century. The main surviving text is in Jer\'onimo de Or\'e's \emph{Rituale seu Manuale Peruanum} (1607). Alfredo Torero wrote his 1965 doctoral thesis on the language. Kallawaya contains Puquina lexical items. A heritage curriculum can only rest on this limited scholarly record.

\section{Primary Educational Resources}
None found.

\section{Linguistic and Reference Materials}
\begin{itemize}
\item Puquina language (Wikipedia): \url{https://en.wikipedia.org/wiki/Puquina_language}
\item Puquina languages (Wikipedia): \url{https://en.wikipedia.org/wiki/Puquina_languages}
\item Leiden University scholarly publication on Puquina (content not inspected beyond search summary): \url{https://scholarlypublications.universiteitleiden.nl/access/item%3A3221120/view}
\item Redalyc article (PDF, content not inspected): \url{https://www.redalyc.org/pdf/3713/371345325004.pdf}
\item Historica (PUCP) article (PDF, content not inspected): \url{https://revistas.pucp.edu.pe/index.php/historica/article/download/31096/27475/125096}
\item Alfredo Torero (Wikipedia biography): \url{https://wikipedia.com/wiki/Alfredo_Torero}
\end{itemize}

\section{Religious and Community Resources}
The 1607 Or\'e \emph{Rituale} is the sole surviving Puquina document, but no online edition was found in search.

\section{Digital and Audio Resources}
None found.

\section{Curriculum Development Context}
No native speakers exist. Content would be limited to history, linguistic heritage, loanwords in southern Peruvian Quechua/Spanish and Kallawaya, and the attested prayer texts.

\end{document}